% Generated. Do not edit.
% Projection scriptural-date-location/evidence, version 1: the Scriptural Date and
% Location page, generated whole from research/chronology.toml, the Scripture
% chronology corpus and the source library by tools/tpt document-types refresh.
% document: liturgy/roman-rite/1962/propers/temporal/55-fifteenth-after-pentecost
\section*{Appendix: Evidence for the Scriptural Dates}
\label{sec:date-location-evidence}
{\small For each passage, as the Scriptural Date and Location table names it: what is dated, how the corpus weighs each answer, the source's own words, and the sources cited. A subject the corpus does not date is reported as such and given no date.\par}
{\small
\par\smallskip\noindent\textit{Offertory: Ps. 39:2--4 (Vg.; Heb. 40)}\par\nobreak
\begin{itemize}[leftmargin=1.2em,itemsep=1pt,topsep=1pt,parsep=0pt]
\item \textbf{Traditional attribution} (disputed). David's reign: B.C.~1055--1015 (``reigned from 1055 to 1015 B.C.''; The Catholic Encyclopedia, Volume 4 (New York, 1908), article 'David, King,' printed p. 642, opening paragraph). Place: Hebron, then Jerusalem (The Holy Bible, Douay-Rheims (Challoner revision), Project Gutenberg eBook 1581: 2 Kings 5:5).
\item \textbf{Historical setting}. No date: the chronology corpus has not yet dated this (research pending).
\item \textbf{Composition} (preferred). Psalms: Before c.~165 B.C. (``before the Maccabean period, around 165 B.C.; no individual psalm can be dated securely''; Catholic critical chronology; New American Bible, Revised Edition --- USCCB web state (2026), Psalms introduction).
\end{itemize}
\par\smallskip\noindent\textit{Introit: Ps. 85:1--3 (Vg.; Heb. 86), Ps. 85:4 (Vg.; Heb. 86)}\par\nobreak
\begin{itemize}[leftmargin=1.2em,itemsep=1pt,topsep=1pt,parsep=0pt]
\item \textbf{Traditional attribution}. As under Offertory: Ps. 39:2--4 (Vg.; Heb. 40).
\item \textbf{Historical setting}. As under Offertory: Ps. 39:2--4 (Vg.; Heb. 40).
\item \textbf{Composition}. As under Offertory: Ps. 39:2--4 (Vg.; Heb. 40).
\end{itemize}
\par\smallskip\noindent\textit{Gradual: Ps. 91:2--3 (Vg.; Heb. 92)}\par\nobreak
\begin{itemize}[leftmargin=1.2em,itemsep=1pt,topsep=1pt,parsep=0pt]
\item \textbf{Historical setting}. As under Offertory: Ps. 39:2--4 (Vg.; Heb. 40).
\item \textbf{Composition}. As under Offertory: Ps. 39:2--4 (Vg.; Heb. 40).
\end{itemize}
\par\smallskip\noindent\textit{Alleluia: Ps. 94:3 (Vg.; Heb. 95)}\par\nobreak
\begin{itemize}[leftmargin=1.2em,itemsep=1pt,topsep=1pt,parsep=0pt]
\item \textbf{Traditional attribution}. As under Offertory: Ps. 39:2--4 (Vg.; Heb. 40).
\item \textbf{Historical setting}. As under Offertory: Ps. 39:2--4 (Vg.; Heb. 40).
\item \textbf{Composition}. As under Offertory: Ps. 39:2--4 (Vg.; Heb. 40).
\end{itemize}
\par\smallskip\noindent\textit{Gospel: Luke 7:11--16}\par\nobreak
\begin{itemize}[leftmargin=1.2em,itemsep=1pt,topsep=1pt,parsep=0pt]
\item \textbf{Narrated event} (preferred). The raising of the widow's son at Naim: A.D.~27 (derived) (``derived A.D. 27 from About Pentecost, A.U.C. 780- Autumn, 780''; The Catholic Encyclopedia, Volume 8 (New York, 1910), article 'Jesus Christ,' III.B.3, printed pp. 378 and 380: chronological calibration and Third and Fifth Journeys).
\item \textbf{Composition} (disputed). The Gospel of St Luke: c.~A.D.~70 (``About the year 70''; The Catholic Encyclopedia, Volume 14 (New York, 1912)); Before the end of the Roman imprisonment, when the Acts was finished (De auctore, de tempore compositionis et de historica veritate Evangeliorum secundum Marcum et secundum Lucam (1912), responsa-i-ix).
\end{itemize}
\par\smallskip\noindent\textit{Communion: John 6:52}\par\nobreak
\begin{itemize}[leftmargin=1.2em,itemsep=1pt,topsep=1pt,parsep=0pt]
\item \textbf{Narrated event} (preferred). The Bread of Life discourse at Capharnaum: A.D.~28 (derived) (``derived Spring A.D. 28 from Spring, A.U.C. 781''; The Catholic Encyclopedia, Volume 8 (New York, 1910), article 'Jesus Christ,' III.B.3, printed pp. 378 and 380: chronological calibration and Third and Fifth Journeys).
\item \textbf{Composition} (disputed). The Gospel of St John: c.~A.D.~90--100 (``from the year 90 to 100 (approximately)''; The Catholic Encyclopedia, Volume 14 (New York, 1912)); A.D.~96--100 (``96 or one of the succeeding years''; The Catholic Encyclopedia, Volume 8 (New York, 1910)). Place: Ephesus traditionally (The Catholic Encyclopedia, Volume 8 (New York, 1910)).
\end{itemize}
\par\smallskip\noindent\textit{Epistle: Gal. 5:25--6:10}\par\nobreak
\begin{itemize}[leftmargin=1.2em,itemsep=1pt,topsep=1pt,parsep=0pt]
\item \textbf{Composition} (answers differ). The Epistle to the Galatians: preferred: A.D.~58 (``anno Christi 58''; Commentaria in omnes D. Pauli epistolas (1614), prooemium; Commentaria in omnes D. Pauli epistolas (1614), epistola-ad-galatas-argumentum); alternative: A.D.~49--50 (``A.D. 49-50''; The Catholic Encyclopedia, Volume 6 (1909)); alternative: c.~A.D.~53--54 (``about A.D. 53, 54''; The Catholic Encyclopedia, Volume 6 (1909)); alternative: A.D.~56 (``(1 and 2 Corinthians; Galatians), 56''; The Catholic Encyclopedia, Volume 11 (New York, 1911)); alternative: A.D.~57--58 (``A.D. 57 or 58''; The Catholic Encyclopedia, Volume 6 (1909)).
\end{itemize}
\par}
